\documentclass[10pt]{article}

\usepackage[utf8]{inputenc}

\usepackage[T1]{fontenc}

\usepackage{amsmath}

\usepackage{amsfonts}

\usepackage{amssymb}

\usepackage[version=4]{mhchem}

\usepackage{stmaryrd}

\usepackage{mathrsfs}

\usepackage{graphicx}

\usepackage[export]{adjustbox}

\graphicspath{ {./images/} }

\usepackage{bbold}



\begin{document}

при $x^{2} \rightarrow 0$ как



\[

\left(-x^{2}+i \varepsilon x_{0}\right)^{1 / 2\left(d_{n}-d_{A}-d_{B}-S_{n}\right)}\left[\ln \left(x^{2} \mu^{2}\right)\right]^{c_{n}},

\]



где $c_{n}$ - числа, к-рые могут быть найдены в рамках теории возмущений. О. р. на световом конусе (3) используется, в частности, для нахождения логарифмич. и степенных эффектов нарушения масштабно-инвариантного поведения структурных функций лептон-адронных глубоко неупругих процессов (см. [3]).



O. р. является эффективным способом вычисления и классификации различных вкладов в физич. амплитуды процессов и находит широкое распространение в приложениях квантовой теории поля. Возможности применения формул (1), (3) в адронной физике связаны с тем, что вид коэффициентных функций $C_{n}^{\mathrm{K}}$, как правило, может быть установлен с помощью теории возмущений, независимо от специфики сильного взаимодействия, после чего сравнение матричных элементов по физич. адронным состояниям от левой и правой частей равенства (1) [или (3)] приводит к соотношениям между физич. амплитудами.



Строгое доказательство О. р. пока существует только в рамках теории возмущений для простых перенормируемых моделей квантовой теории поля (см. [4]).



Mum.: [1] Wils on K.G., "Phys. Rev.», 1969, v. 179, p. 1499; [2] Shifman M.A., Vainshte in A. I., Zakharov V.I., "Nucl. Phys. B», 1979, v. 147, p. 385; [3] Gross D.J., Wilczek F., «Phys. Rev. D», 1974, v. 9, p. 980; [4] За в ья л о в О. И., Перенормированные диаграммы Фейнмана, М., $1979 . \quad$ Л. Н. Липатов. OПEPATOPHOЗНАЧНАЯ ФУНКЦИЯ - элемент множества $\mathscr{F}(X, \mathscr{L}(Y, Z))$, где $\mathscr{F}(X, \mathscr{L})-$ множество функций, заданных на множестве $X$, со значением во множестве $\mathscr{L}$, а $\mathscr{L}(Y, Z)$ - множество линейных операторов, отображающих



\begin{center}

\includegraphics[max width=\textwidth]{2023_07_08_4ae0c9e27b25a37b707fg-1}

\end{center}



Пусть $X$ - банахово пространство, и пусть $D$ - область комплексной плоскости (связное открытое множество). Функция $x(\cdot)$, определенная на $D$, со значением в $X$ называется с ильно голоморфной в $z_{0}$, если сушествует сильный предел отношения



\[

\frac{x\left(z_{0}+h\right)-x\left(z_{0}\right)}{h}

\]



при $h$, стремящемся к нулю в $\mathbb{C}$.



Сильно голоморфная функция в каждой точке области $D$ комплексной плоскости со значением в банаховом пространстве ядерных операторов называется О. Ф., аналитиче$\begin{array}{ll}\text { ской в ядерной норме. } & \text { В. И. Скрипник. }\end{array}$ OПEPATOPHOЗНАЧНЫЙ СИМВОЛ Д ифф е ре н циального или псе вдодифференциального о п е р а то р а - функция со значениями в некоторой операторной алгебре, характеризующая данный оператор, структура которого оказывается более простой, чем у исходного оператора. В задачах с малым параметром $h$ дифференциальные или псевдодифференциальные операторы с O.с. имеют вид: $L(-i h \nabla, x, h)$, где $L(p, x, h)-$ О. с.



П р и м е р: символ $L(p, x)(m \times m)$-матричного $h$-дифференциального оператора - $(m \times m)$-матричная функция (оператор, действуюший в конечномерном пространстве).



Другой п р и м е р дают уравнения, в к-рых производные по части переменных $\left(x_{1}, x_{2}, \ldots, x_{n}\right)$ содержат параметр, а по другим переменным $\left(y_{1}, y_{2}, \ldots, y_{l}\right)$ его не содержат. Тогда $L(p, x, h)$ - при каждых $(p, x)$ - дифференциальный оператор по переменным $y$. Напр., символ $L(p, x)$ уравнения Шредингера для взаимодействуюших легких и тяжелых частиц имеет вид:



\[

L(p, x)=p^{2}-\frac{\partial^{2}}{\partial y_{1}^{2}}-\frac{\partial^{2}}{\partial y_{2}^{2}}-\ldots-\frac{\partial^{2}}{\partial y_{l}^{2}}+v(x, y) .

\]



Лum.: [1] М ас ло в В. П., Асимптотические методы и теория возмущений, М., 1988; [2] е го же, Теория возмущений и асимптотические методы, М., 1965.



С. Ю. Доброхотов. OREPATOPHEIЙ METOД - см. Равномерно пригодные приближения в квантовой механике.



OREPATOPH в к в а н т о в о й те о р и и - понятие, широко используемое в математическом аппарате квантовой механики и квантовой теории поля. О. служат для сопоставления с определенной волновой функцией (или вектором состояния) $\Psi$ другой определенной функции (вектора) $\Psi^{\prime}$. Соотношение между $\psi$ и $\psi^{\prime}$ записывается в виде $\psi^{\prime}=\hat{L} \psi$, где $\hat{L}-$ O. В квантовой механике физич. величинам $L$ (координате, импульсу, энергии и др.) ставятся в соответствие О. $\hat{L}$ (оператор координаты, импульса и т. д.), действующие на $\psi$. Простейшие виды О., действуюшие на волновую функцию $\Psi(x)$ (где $x-$ координата частицы),- О. умножения (напр., О. координаты $\hat{x}, \hat{x} \psi=x \psi$ ) и О. дифференцирования (напр., О. $\hat{p}_{x}$ проекции импульса на ось $x, \hat{p}_{x} \psi=-i \hbar \partial \psi / \partial x$ ). Если $\Psi$ - вектор, компоненты к-рого можно представить в виде столбца чисел, то О. представляет собой квадратную матрицу.



В квантовой механике в основном используются линейные О., к-рые обладают следующим свойством: если



\[

\hat{L} \psi_{1}=\psi_{1}^{\prime} \text { и } \quad \hat{L} \psi_{2}=\psi_{2}^{\prime}

\]



To



\[

\hat{L}\left(c_{1} \Psi_{1}+c_{2} \psi_{2}\right)=c_{1} \psi_{1}^{\prime}+c_{2} \psi_{2}^{\prime}

\]



где $c_{1}$ и $c_{2}-$ комплексные числа.



Свойства O. $\hat{L}$ определяются уравнением



\[

\hat{L} \Psi_{n}=\lambda_{n} \Psi_{n}

\]



где $\lambda_{n}$ - числа. Решения этого уравнения $\psi_{n}$ называются собственными функциями (собственными векторами) O. $\hat{L}$. Собственные волновые функции (собственные векторы состояния) описывают в квантовой механике такие состояния, в к-рых физич. величина $L$ (соответствующая О. $\hat{L})$ имеет определенное значение $\lambda_{n}$. Числа $\lambda_{n}$ называются собственными значениями О. $\hat{L}$, а их совокупность - спектром оператора. Спектр может быть непрерывным или дискретным; в первом случае уравнение, определяющее $\psi_{n}$, имеет решение при любом значении $\lambda_{n}$ (в определенной области), во втором - решения существуют только при определенных дискретных значениях $\lambda_{n}$. Спектр О. может быть и смешанным: частично непрерывным, частично дискретным. Напр., О. координаты и импульса имеют непрерывный спектр, а О. энергии в зависимости от характера действуюших в системе сил - непрерывный, дискретный или смешанный спектр. Дискретные собственные значения О. энергии называются у ровня м и эне р г и и.



Собственные функции и собственные значения О. физич. величин должны удовлетворять определенным требованиям. Так как непосредственно измеряемые физич. величины всегда принимают вещественные значения, то соответствующие квантовомеханич. О. должны иметь вещественные собственные значения. Поскольку при измерении физич. величины в любом состоянии $\psi$ должно получаться одно из возможных собственных значений О. этой величины, необходимо, чтобы произвольная волновая функция (вектор состояния) могла быть представлена в виде линейной комбинации собственных функций (векторов) $\psi_{n}$ О. этой физич. величины; другими словами, совокупность собственных функций (векторов) должна представлять полную систему. Этими свойствами обладают собственные функции и собственные значения так наз. самосопряженных, или эрмитовых, $\mathrm{O}$.



С О. можно производить алгебраич. действия. В частности, под п ро и з вед е ние м О. $\hat{L}_{1}$ и $\hat{L}_{2}$ понимается такой О. $\hat{L}=\hat{L}_{1} \hat{L}_{2}$, действие к-рого на $\psi$ дает $\hat{L} \psi=\psi^{\prime \prime}$, если $\hat{L}_{2} \Psi=\psi^{\prime}$ и $\hat{L}_{1} \Psi^{\prime}=\psi^{\prime \prime}$. Произведение О. в общем случае зависит от порядка сомножителей, то есть $\hat{L}_{1} \hat{L}_{2} \neq \hat{L}_{2} \hat{L}_{1}$. Этим алгебра О. отличается от обычной алгебры чисел. Возможность пере-





\end{document}